\begin{ex}
Viết thương của phép chia $2 : 3$ dưới dạng phân số.
\shortans{2/3}
\loigiai{
$2 : 3 = \frac{2}{3}$.
}
\end{ex}

\begin{ex}
Viết thương của phép chia $4 : 5$ dưới dạng phân số.
\shortans{4/5}
\loigiai{
$4 : 5 = \frac{4}{5}$.
}
\end{ex}

\begin{ex}
Viết thương của phép chia $7 : 9$ dưới dạng phân số.
\shortans{7/9}
\loigiai{
$7 : 9 = \frac{7}{9}$.
}
\end{ex}

\begin{ex}
Viết thương của phép chia $9 : 4$ dưới dạng phân số.
\shortans{9/4}
\loigiai{
$9 : 4 = \frac{9}{4}$.
}
\end{ex}

\begin{ex}
Viết số tự nhiên 5 dưới dạng phân số có mẫu số bằng 1.
\shortans{5/1}
\loigiai{
$5 = \frac{5}{1}$.
}
\end{ex}

\begin{ex}
Viết số tự nhiên 12 dưới dạng phân số có mẫu số bằng 1.
\shortans{12/1}
\loigiai{
$12 = \frac{12}{1}$.
}
\end{ex}

\begin{ex}
Viết số tự nhiên 1 dưới dạng phân số có tử số bằng 8.
\shortans{8/8}
\loigiai{
$1 = \frac{8}{8}$.
}
\end{ex}

\begin{ex}
Viết số tự nhiên 0 dưới dạng phân số có mẫu số bằng 6.
\shortans{0/6}
\loigiai{
$0 = \frac{0}{6}$.
}
\end{ex}

\begin{ex}
Phân số $\frac{12}{3}$ có giá trị bằng số tự nhiên nào?
\shortans{4}
\loigiai{
$12 : 3 = 4$.
}
\end{ex}

\begin{ex}
Phân số $\frac{35}{5}$ có giá trị bằng số tự nhiên nào?
\shortans{7}
\loigiai{
$35 : 5 = 7$.
}
\end{ex}

\begin{ex}
Phân số $\frac{48}{6}$ có giá trị bằng số tự nhiên nào?
\shortans{8}
\loigiai{
$48 : 6 = 8$.
}
\end{ex}

\begin{ex}
Phân số $\frac{72}{8}$ có giá trị bằng số tự nhiên nào?
\shortans{9}
\loigiai{
$72 : 8 = 9$.
}
\end{ex}

\begin{ex}
Có 3 cái bánh chia đều cho 4 bạn. Mỗi bạn được bao nhiêu phần cái bánh?
\shortans{3/4}
\loigiai{
Mỗi bạn được $\frac{3}{4}$ cái bánh.
}
\end{ex}

\begin{ex}
Có 5 lít nước chia đều vào 6 chai. Mỗi chai chứa bao nhiêu lít nước?
\shortans{5/6}
\loigiai{
Mỗi chai chứa $\frac{5}{6}$ lít.
}
\end{ex}

\begin{ex}
Có 7 kg đường chia đều vào 8 túi. Mỗi túi chứa bao nhiêu ki-lô-gam đường?
\shortans{7/8}
\loigiai{
Mỗi túi chứa $\frac{7}{8}$ kg.
}
\end{ex}

\begin{ex}
Chia đều 9 m vải thành 5 đoạn bằng nhau. Mỗi đoạn dài bao nhiêu mét?
\shortans{9/5}
\loigiai{
Mỗi đoạn dài $\frac{9}{5}$ m.
}
\end{ex}

\begin{ex}
Trong các phân số $\frac{3}{4}, \frac{5}{5}, \frac{7}{6}, \frac{8}{9}$, phân số nào bằng 1?
\shortans{5/5}
\loigiai{
Phân số $\frac{5}{5} = 1$.
}
\end{ex}

\begin{ex}
Trong các phân số trên, có bao nhiêu phân số bé hơn 1?
\shortans{2}
\loigiai{
$\frac{3}{4}$ và $\frac{8}{9}$ (có 2 phân số).
}
\end{ex}

\begin{ex}
Trong các phân số trên, có bao nhiêu phân số lớn hơn 1?
\shortans{1}
\loigiai{
$\frac{7}{6}$ (có 1 phân số).
}
\end{ex}

\begin{ex}
Tìm số tự nhiên $x$ biết: $\frac{x}{7} = 1$.
\shortans{7}
\loigiai{
$x = 7$.
}
\end{ex}

\begin{ex}
Tìm số tự nhiên $x$ biết: $\frac{x}{5} = 4$.
\shortans{20}
\loigiai{
$x = 4 \times 5 = 20$.
}
\end{ex}

\begin{ex}
Tìm số tự nhiên $y$ biết: $\frac{18}{y} = 6$.
\shortans{3}
\loigiai{
$y = 18 : 6 = 3$.
}
\end{ex}

\begin{ex}
Mọi số tự nhiên đều có thể viết dưới dạng phân số có mẫu số bằng mấy?
\shortans{1}
\loigiai{
Mẫu số bằng 1.
}
\end{ex}

\begin{ex}
Phép chia $0 : 9$ viết dưới dạng phân số là bao nhiêu?
\shortans{0/9}
\loigiai{
$\frac{0}{9}$.
}
\end{ex}

\begin{ex}
Thương của phép chia số tự nhiên cho số tự nhiên khác 0 luôn viết được dưới dạng gì? (Phân số ghi 1, số thập phân ghi 2).
\shortans{1}
\loigiai{
Phân số (1).
}
\end{ex}
